% ======================= 摘要专用页 =======================
\begin{center}
  {\heiti\zihao{3} 微网与外部电网电力调控策略的优化模型}\\[6pt]
  {\heiti\zihao{4} 摘\quad 要}
\end{center}

\begingroup
\zihao{-4}

本文研究微网在电价、负载与光伏出力三重信息条件下的购电与储能调控问题，将四个问题\textbf{统一归结为含储能的单母线能量流线性规划}：以时段末储电量 $E_t$、计划购电量 $g_t$ 与净缺口 $n_t$ 为核心量，按"确定信息—预测不确定—多阶段更新—电价波动"递进引入不确定性，全部模型均为线性规划，可用 HiGHS 精确求解。

\textbf{针对问题1}，建立 144 个时段的确定性线性规划，约束含母线功率平衡、储能动态、容量与功率限制及日周期条件 $E_{144}=E_0$。由电量守恒导出全天购电量分解式 $\sum_tg_t=\text{净缺口}+\text{弃光}+(1/\eta_c-\eta_d)A$；由 KKT 条件证明最优策略为\textbf{影子价格阈值策略}（$p_t<0.9\lambda$ 充电、$p_t>1.111\lambda$ 放电，套利门槛价比 $1.235$）。求得全天购电费 \textbf{35101.57 元}、全天购电量 59482.70 \si{kWh}，较无储能基线节省 12950.48 元（\textbf{27.0\%}），0:00 与 24:00 储电量均为 6000 \si{kWh}。

\textbf{针对问题2}，建立"日前预测—安全裕度—实际结算"的两阶段模型：经滚动精度评价与决策导向标定，负载取 7 日滑动平均（RMSE 955.9 \si{kW}）、光伏取 EWMA(0.7,14)（MAE 150.1 \si{kW}）；安全裕度以 1 月数据按总费用最小标定为 $(\alpha,\beta)=(0.15,0)$；报童模型给出最优覆盖分位点 $c_u/(c_u+c_o)\in[0.80,0.96]$，与标定结果互相印证。全年总费用 \textbf{1664.80 万元}（计划 1534.78、紧急 130.02），紧急购电量仅占负载的 \textbf{0.746\%}，安全裕度节省 \textbf{496.37 万元（22.97\%）}。

\textbf{针对问题3}，将"调低罚 0.5 倍、调高收 1.5 倍"的三档价格化简为对称罚金 $\Phi_t=p_t\tilde g_t+0.5p_t|\tilde g_t-\hat g_t|$，据此证明调整机制下的最优计划分位点为\textbf{中位数}，并给出更新价值分解 $V_\tau\approx0.5p\cdot\Delta\mathrm{MAD}+\Delta(\text{紧急费})$。滚动时域重优化下题目方案 P3 的总费用为 \textbf{2136.95 万元}。实测 12:00 更新的边际价值为正（$5.25$ 万元），6:00 与 18:00 为负，逐小时更新亦无增益（且 19:00--24:00 四个发布时刻的预报完全相同）；信息价值分解显示完美光伏信息价值 $200.25$ 万元，仅为完美负载信息价值 $520.39$ 万元的 $38\%$。故\textbf{本题数据条件下不需要引入其他时刻的预报}，制约收益的是负载侧不确定性。

\textbf{针对问题4}，以 7 日滑动平均预测电价、按实际电价结算重算问题2 与问题3，总费用分别为 \textbf{2034.59 万元}与 \textbf{2224.44 万元}（较固定电价上升 $22.2\%$ 与 $4.1\%$）。由于电价与负载的相关系数为 $0.508$，本文证明此时最优计划分位点\textbf{高于中位数}，与鲁棒模型"取价格上界"的结论一致。

\textbf{模型检验}方面，用独立线性规划（母线侧变量布局）与独立动态规划交叉验算确定性核心（偏差 $7\times10^{-8}$ 与 $0.276\%$），用独立结算公式复算全部费用并核对 5 个结果文件的结构与守恒关系，共 \textbf{40 项检查全部通过}；灵敏度分析表明充放电效率（约 $\pm4\%$）与储电量上下限（约 $2.6\%$）是主要影响参数，而日周期约束下初始电量对费用几乎无影响。全文以线性规划为统一内核，用影子价格、报童分位点与信息价值分解三条理论线索串联四问，使"策略为何如此"与"结果是多少"相互印证。

\vspace{6pt}
\noindent{\heiti 关键词：}\quad 微网购电优化；储能套利；安全裕度；滚动时域优化；信息价值；线性规划
\endgroup
